\documentclass[sigconf,anonymous,review]{acmart}

\usepackage{amsmath}
\usepackage{booktabs}
\usepackage{multirow}
\usepackage{xcolor}

\settopmatter{printacmref=false}
\renewcommand\footnotetextcopyrightpermission[1]{}

% Draft-only marker. Replace every instance with an author-verified citation.
\newcommand{\citetodo}[1]{\textcolor{red}{[CITATION: #1]}}
\newcommand{\system}{HCPM}

\title{HCPM: Discovering Latent Habits for Controlled Proactive Memory in LLM Agents}

\author{Anonymous Author(s)}
\affiliation{\institution{Anonymous Institution}}

\begin{document}

\begin{abstract}
Long-term memory allows large language model agents to retain information beyond a single interaction, but retrieval alone does not determine when remembered information should interrupt an ongoing task.
This distinction is important for web-facing and tool-using agents: users often cannot fully articulate their recurring workflows, while an agent that surfaces every related memory becomes distracting.
We formulate proactive memory delivery as a turn-level decision problem and introduce \system{}, a middleware that discovers latent habits from interaction traces and controls whether they should be surfaced at the current turn.
\system{} consolidates repeated event--consequence patterns into typed habits, generalizes variable event forms, generates candidates from multiple trigger families, and applies scoring, phase-aware gating, and per-habit cooldown before producing a structured Push or no-Push decision.
Across six domain adaptations covering code, smart home, shopping, email, calendar, and project management, \system{} reaches a macro F1 of 0.653, compared with 0.392 for an oracle that selects the strongest baseline independently in each domain.
It also reduces false positives from 604 to 176 and proactive pushes from 822 to 408 relative to those selected baselines.
Ablations show that generalized habit matching is the broadest contributor, while phase gating has a large, repeatable Code-specific effect.
Cooldown exposes a measurable recall--intervention trade-off rather than universally maximizing F1.
These results support a modular view of proactive memory: useful systems must discover unspoken regularities and separately control when acting on them is warranted.
\end{abstract}

\keywords{LLM agents, long-term memory, habit discovery, proactive assistance, intervention control}

\maketitle

\section{Introduction}

Large language model (LLM) agents increasingly operate over web applications, workplace tools, software repositories, and personal services.
In these settings, an agent may interact with the same user and environment over many sessions.
Persistent memory systems address this setting by storing previous events, retrieving relevant records, summarizing interaction history, or organizing memory into structured stores \citetodo{Generative Agents; MemGPT; MemoryBank; Mem0; A-Mem; HippoRAG}.
These mechanisms answer an important question: \emph{what past information is relevant now?}
They do not, by themselves, answer a second question: \emph{should that information be surfaced at the current turn?}

The second question becomes central when memory is proactive.
A reactive system waits for a query; a proactive system may interrupt an ongoing workflow.
The cost of an incorrect decision is therefore asymmetric.
Missing a useful reminder loses an opportunity, while an unnecessary reminder consumes attention, may expose an incorrect inference about the user, and can steer an agent toward an action that was never intended.
Relevance is not enough: a memory can be related to the current task but premature, already satisfied, inconsistent with an explicit exception, or repeated too soon after an earlier intervention.

Recent systems have begun to study proactive demand detection, reasoning-integrated memory navigation, and selective reminder injection \citetodo{PASK; MemCog; Remember When It Matters}.
We focus on a complementary gap.
Users frequently cannot enumerate their own recurring workflows or express them as stable rules.
For example, a developer may regularly commit only after tests pass, a shopper may repeatedly avoid a class of ingredients, or a calendar user may follow a scheduling preference without ever declaring it as a memory.
An agent should be able to infer such regularities from interaction traces.
At the same time, the inferred habit must not become an unconditional notification rule.

We introduce \system{} (Habit-Consolidated Proactive Memory), a lightweight middleware for discovering latent habits and controlling their delivery.
\system{} observes structured user and agent events, accumulates repeated evidence, and consolidates trigger--context--consequence regularities into typed habit records.
It generalizes semantically equivalent event forms, generates proactive candidates through several trigger families, and exposes an auditable controller that scores candidates, checks workflow phase, optionally consults an LLM for high-risk cases, and enforces cooldown.
Its evaluated output is a structured Push or no-Push decision; optional natural-language realization occurs only after this decision and is outside the core mechanism.

This formulation is relevant to web search and data mining in two ways.
First, web agents and intelligent assistants must mine longitudinal interaction traces rather than treat each request independently.
Second, the utility of retrieved behavioral evidence depends on decision timing: the system must rank candidate memories and decide whether any candidate should enter the current agent context.
The resulting problem joins behavioral pattern mining, retrieval, and selective intervention for agentic systems.

We make four contributions:
\begin{itemize}
    \item We formulate proactive memory as a turn-level decision problem and construct six domain adaptations from repository interactions and public WorkBench, $\tau$-bench Retail, and SmartBench sources. These are adaptations, not official-protocol benchmark scores.
    \item We present \system{}, which discovers habits that need not be explicitly declared and couples generalized habit matching with an auditable intervention controller.
    \item We provide six-domain effectiveness, ablation, statistical, repeatability, and error evidence. \system{} improves macro F1 from 0.392 for the strongest per-domain baselines to 0.653, while its mechanisms activate under different domain structures.
    \item We test targeted mechanism effects across four model configurations and report absolute call, token, retry, and latency composition without claiming matched-baseline efficiency superiority.
\end{itemize}

\section{Problem Formulation}

\subsection{Events, Habits, and Decisions}

Let a session be an ordered stream of events $E=(e_1,\ldots,e_T)$.
Each event is represented as
\begin{equation}
e_t=(\tau_t,a_t,v_t,o_t,m_t),
\end{equation}
where $\tau_t$ is time, $a_t$ is the actor, $v_t$ is an action verb, $o_t$ is an object or entity, and $m_t$ contains domain metadata.
Structured tool calls are mapped deterministically; free-form language can be parsed into the same schema by an LLM.
The benchmark uses pre-structured events so event parsing is not part of the reported decision accuracy.

A habit $h$ represents a recurring behavioral regularity:
\begin{equation}
h=(g_h,c_h,y_h,q_h,f_h,r_h,z_h,V_h).
\end{equation}
Here, $g_h$ is a trigger action--entity pattern, $c_h$ is trigger context, $y_h$ is a distribution over typical consequences, $q_h$ is confidence, $f_h$ and $r_h$ encode frequency and recency, $z_h$ is lifecycle state, and $V_h$ is a set of generalized variants.
Habit types include workflow, temporal, emotional, and associative patterns.

At each turn, the system observes the current event and memory state $M_t$, and outputs
\begin{equation}
\hat{y}_t = \pi(e_{\leq t},M_t) \in \{0,1\},
\end{equation}
where $1$ denotes a structured Push decision.
When $\hat{y}_t=1$, the decision also identifies its trigger, supporting habit, candidate score, and control trace.
Natural-language wording is a downstream host-system concern and does not change $\hat{y}_t$.

\subsection{Evaluation Objective}

Each evaluation turn has a ground-truth label $y_t$ indicating whether an intervention is useful at that stage.
We report precision, recall, F1, accuracy, false positives (FP), and Push count.
F1 captures the balance between recovering useful interventions and avoiding false ones, while FP and Push count expose intervention burden directly.
No single metric fully represents user utility: the relative cost of one missed intervention and one unnecessary interruption is deployment dependent.
We therefore treat threshold and cooldown as an operating policy rather than tune them solely to maximize test F1.

\section{HCPM Framework}

\subsection{Overview}

Figure~\ref{fig:architecture} summarizes the pipeline.
\system{} separates two responsibilities.
The \emph{discovery layer} turns repeated interaction evidence into typed, persistent habits.
The \emph{control layer} decides whether a habit-derived candidate is actionable at the current turn.
This separation prevents ``memory found'' from becoming synonymous with ``notification sent.''

\begin{figure*}[t]
    \centering
    \fbox{\parbox{0.96\textwidth}{\centering
    \textbf{Figure 1 placeholder.} Interaction events $\rightarrow$ event counter and batch evidence $\rightarrow$ generalized habit discovery $\rightarrow$ typed Habit Store $\rightarrow$ multi-source candidates $\rightarrow$ score, phase/Think, threshold, and cooldown $\rightarrow$ structured Push/no-Push. Optional language realization is outside the evaluated decision boundary.}}
    \caption{\system{} separates discovery of unspoken habits from control of proactive intervention. The final production figure will be loaded from \texttt{paper/assets/fig1\_architecture/}.}
    \label{fig:architecture}
\end{figure*}

\subsection{Event Accumulation and Habit Discovery}

An event counter maintains frequencies of $(o,v)$ pairs and a sliding buffer of recent event summaries.
Every discovery batch filters out conversational, terminal, and verification actions that should not independently define a habit.
Candidates must have at least $n_{\min}=3$ supporting events in the default configuration.
For coding events, deterministic preprocessing merges repeated files with the same extension (e.g., several Python edits can become \texttt{edit *.py}); domains with normalized entities retain their verb--object keys.

The remaining evidence is passed to an LLM with a constrained schema.
The output proposes a habit type, trigger, trigger context, consequence, and variants.
The parser removes reasoning blocks, scans for a schema-valid JSON payload, and retries malformed outputs before applying domain-specific fallbacks.
Thus, the LLM proposes structured regularities, while deterministic checks govern admissible evidence and the representation written to memory.
Users need not state ``this is my habit,'' but the system still requires repeated observational evidence and does not claim zero-history personalization.

\subsection{Habit Confidence and Lifecycle}

Each record tracks frequency, last observation, consequence statistics, confidence, and lifecycle state.
Let
\begin{align}
F(h) &= \min(f_h/f_{\max},1),\\
R(h) &= \exp(-\lambda \Delta_h),
\end{align}
where $\Delta_h$ is days since the last observation.
Consistency $C(h)$ decreases when a trigger has many distinct consequences.
The implementation computes
\begin{equation}
q_h=\sigma\!\left(6\left[w_fF(h)+w_rR(h)+w_cC(h)-0.5\right]\right),
\end{equation}
with $(w_f,w_r,w_c)=(0.5,0.3,0.2)$ and $\lambda=0.1$.
Records transition among active, fading, forgotten, and abandoned states; repeated evidence can reinforce a fading record.
Only the Habit Store requires persistence, while counters and queues can be reconstructed from the event stream.

\subsection{Generalized Habit Matching}

Exact matching is brittle when behavior is stable but surface forms vary.
\system{} therefore stores generalized trigger patterns and variants.
A coding habit may match a glob pattern, while preference-oriented domains can map an observed conflict key to a normalized habit entity.
At runtime, Variant Match selects the highest-confidence supported variant and emits a candidate describing the typical consequence.
The Exact-only ablation retains habit discovery but disables generalized runtime matching, isolating the value of abstraction beyond literal event equality.

\subsection{Multi-Source Candidate Generation}

The controller supports seven trigger families:
\begin{itemize}
    \item \textbf{Semantic}: retrieves a related record from an external vector store.
    \item \textbf{Pattern}: fires when an expected consequence remains unfulfilled past its time window.
    \item \textbf{Anomaly}: detects a structurally related action that deviates from a high-confidence habit.
    \item \textbf{Variant Match}: matches the current action to a discovered exact or generalized habit variant.
    \item \textbf{Temporal}: activates a typed temporal habit at a compatible time.
    \item \textbf{Emotional}: matches stored emotional cues or keywords.
    \item \textbf{Associative}: surfaces a habit through an associated concept.
\end{itemize}

Each candidate $c$ has relevance $r_c$, urgency $u_c$, confidence $q_c$, and novelty $n_c$.
The score is
\begin{equation}
s(c)=0.35r_c+0.25u_c+0.25q_c+0.15n_c.
\end{equation}
The highest-scoring candidate proceeds only if $s(c)>\theta_{\mathrm{push}}$, with $\theta_{\mathrm{push}}=0.6$ by default.

\subsection{Phase-Aware and Confidence-Aware Control}

Related behavior is not necessarily actionable behavior.
For workflow domains with observable execution stages, \system{} derives a compact phase from recent events.
In Code, for example, current test status and edit/run/debug activity distinguish early work, implementation, debugging, documentation work, and wrap-up.
A matching habit can be suppressed during early work or active debugging and allowed after stable verification.
For non-coding domains where those signals do not apply, the phase module returns an active state rather than imposing coding-specific rules.

Pattern and Anomaly candidates can additionally pass through an optional LLM Think check.
The checker receives current phase, recent actions, test state, and the candidate, then returns PUSH or SUPPRESS with calibrated confidence.
Failures fall back to the structured decision path and are recorded.
As Section~\ref{sec:ablation} shows, this checker has weak marginal evidence in the current experiments; it is retained as an optional safeguard rather than presented as the central contribution.

\subsection{Cooldown and Feedback Control}

After thresholding, a candidate is checked against a cooldown key containing trigger, source, and habit identifier.
This design prevents one recently surfaced habit from blocking a different habit.
The frozen experiments use a three-event cooldown, reset at session boundaries.
After an observed response to a Push, a global feedback cooldown can also suppress immediate follow-up notifications, whether the suggestion was used or ignored.

Cooldown determines an intervention operating point.
It is not a categorical user-interface label and is not expected to maximize F1 in every domain.
The complete decision trace stores the candidate, phase, Think result, cooldown distance, suppression reason, and final outcome, allowing an evaluator to distinguish absent evidence from deliberate suppression.

\subsection{Deployment Boundary}

The structured decision is middleware output.
An agent host may inject it directly, transform it into a notification, request confirmation, or suppress it under an application policy.
Optional LLM rewriting occurs after selection and cooldown and therefore does not alter decision F1.
We report this cost separately in Section~\ref{sec:efficiency}.

\section{Experimental Setup}

\subsection{Research Questions}

We study four questions: (RQ1) Does \system{} improve proactive decision quality over common memory baselines? (RQ2) Which mechanisms contribute under different domain structures? (RQ3) Are the effects stable across repetitions, protocols, and model configurations? (RQ4) What is the measured call and token composition of the framework?

\subsection{Domains and Data Construction}

Table~\ref{tab:data} summarizes the evaluation data.
Code uses repository-interaction scenarios.
Email, Calendar, and Project Management are adaptations derived from WorkBench \citetodo{WorkBench}; Shopping is adapted from $\tau$-bench Retail \citetodo{tau-bench}; SmartHome is adapted from the public SmartBench repository \citetodo{SmartBench GitHub repository}.
We transform these sources into longitudinal event sessions with turn-level labels for whether a proactive memory should be delivered.
The reported numbers are not directly comparable to official, unmodified benchmark leaderboards.

\begin{table}[t]
\centering
\caption{Dataset statistics. C/H denote clean/hard subsets. Warm-up events are excluded from evaluated turns.}
\label{tab:data}
\small
\begin{tabular}{lrrrr}
\toprule
Domain & C sess. & H sess. & Turns & Warm-up \\
\midrule
Code & 20 & 12 & 448 & 34 \\
SmartHome & 20 & 12 & 440 & 50 \\
Shopping & 20 & 12 & 495 & 54 \\
Email & 25 & 12 & 684 & 54 \\
Calendar & 25 & 12 & 684 & 54 \\
Project Mgmt. & 25 & 12 & 684 & 54 \\
\bottomrule
\end{tabular}
\end{table}

Clean sessions contain clearer evidence and expected intervention stages.
Hard sessions add premature cues, implicit positives, aligned actions, and explicit exceptions.
The main \texttt{all} condition is a continual stateful clean-to-hard run: the Habit Store, counter state, and learned variants persist across the subset boundary, while session-local cooldown is reset.
We separately audited split-reset and reversed-order protocols rather than treating \texttt{all} as a simple arithmetic union.

\subsection{Baselines}

All methods receive the same current event and output Push/no-Push under a common runner.
\textbf{No Memory} uses only the current event and short local context.
\textbf{RAG-only} stores raw historical events and retrieves similar observations.
\textbf{Summary-Reflection} compresses observations into a long-term natural-language summary and reflections.
\textbf{Generative Memory} follows a relevance--recency--importance retrieval scheme inspired by Generative Agents \citetodo{Generative Agents}.
The baselines intentionally exclude \system{}'s typed Habit Store, generalized preference abstraction, and consolidation mechanism.
Warm-up and a per-memory cooldown of three events are applied where supported to avoid giving \system{} an intervention-frequency advantage by protocol alone.

\subsection{Metrics and Statistical Tests}

We report precision, recall, F1, accuracy, TP, FP, FN, TN, and Push count from complete decision traces.
For ablations, 95\% confidence intervals use 10,000 paired, stratified session-level bootstrap samples, preserving within-session dependence and clean/hard stratum sizes.
We also use two-sided exact McNemar tests on turn-level correctness, with Holm family-wise correction across 14 comparisons.
The two analyses answer different questions: bootstrap intervals test F1 effects, whereas McNemar tests symmetry in correctness changes.

\subsection{Implementation and Reproducibility}

The main matrix uses DeepSeek-v4-Flash with temperature zero, domain-specific warm-up, LLM Think enabled, and cooldown three.
The package requires Python 3.10 or newer; later audit scripts used Python 3.12.13.
Every evaluated turn writes a JSONL trace containing trigger hits, candidates, score components, phase, Think output, cooldown state, ground truth, and TP/FP/FN/TN classification.
The final audit validates 90/90 main-result logs and 36,779/36,779 manifest-backed trace rows across 65 runs.
API keys and endpoint credentials are excluded from the artifact.

\section{Results and Analysis}

\subsection{Overall Performance}

Table~\ref{tab:main} reports F1 on the continual \texttt{all} protocol.
\system{} has the highest F1 in all six domains.
Its macro F1 is 0.653, compared with 0.388 for RAG-only and 0.392 for an oracle that chooses the strongest baseline separately in each domain.

\begin{table*}[t]
\centering
\caption{F1 on six continual \texttt{all} adaptations under DeepSeek-v4-Flash. Macro is the unweighted mean over domains. The oracle baseline macro (0.392) selects the strongest baseline independently per domain and is not a deployable single method.}
\label{tab:main}
\small
\begin{tabular}{lrrrrrrr}
\toprule
Method & Code & SmartHome & Shopping & Email & Calendar & Project Mgmt. & Macro \\
\midrule
No Memory & 0.000 & 0.087 & 0.355 & 0.239 & 0.295 & 0.034 & 0.168 \\
RAG-only & 0.247 & 0.380 & 0.439 & 0.494 & 0.381 & 0.387 & 0.388 \\
Summary-Reflection & 0.176 & 0.314 & 0.361 & 0.370 & 0.408 & 0.315 & 0.324 \\
Generative Memory & 0.197 & 0.360 & 0.427 & 0.386 & 0.393 & 0.335 & 0.350 \\
\textbf{HCPM} & \textbf{0.667} & \textbf{0.647} & \textbf{0.600} & \textbf{0.683} & \textbf{0.672} & \textbf{0.652} & \textbf{0.653} \\
\bottomrule
\end{tabular}
\end{table*}

Compared with the strongest baseline in each domain, the F1 gains are 0.420 (Code), 0.267 (SmartHome), 0.161 (Shopping), 0.188 (Email), 0.264 (Calendar), and 0.265 (Project Management).
The selected baselines produce 604 false positives and 822 pushes, whereas \system{} produces 176 false positives and 408 pushes, reductions of 70.9\% and 50.4\%, respectively.

The aggregate does not imply uniform metric dominance.
In Shopping, recall-heavy RAG-only reaches 1.000 recall but produces 133 false positives; \system{} reaches 0.635 recall with 25 false positives.
In Calendar, \system{} improves F1 and recall but produces 33 false positives, seven more than the best-F1 Summary-Reflection baseline.
These cases motivate reporting intervention count alongside F1.

\subsection{Mechanism Ablations}
\label{sec:ablation}

Table~\ref{tab:ablation} shows a modular rather than universally additive system.
Exact-only retains discovered habits but removes generalized runtime matching.
Removing the phase gate changes Code substantially but is prediction-identical in four other tested domains.

\begin{table*}[t]
\centering
\caption{Ablation F1; parentheses are changes from Full. A dash denotes not tested.}
\label{tab:ablation}
\small
\begin{tabular}{lrrrrr}
\toprule
Domain & Full & Exact-only & w/o Phase & w/o Think & w/o Cooldown \\
\midrule
Code & \textbf{0.667} & 0.667 (0.000) & 0.462 ($-$0.205) & 0.642 ($-$0.024) & -- \\
Shopping & 0.600 & 0.429 ($-$0.171) & 0.600 (0.000) & 0.600 (0.000) & \textbf{0.701} (+0.101) \\
Calendar & \textbf{0.672} & 0.133 ($-$0.539) & -- & -- & -- \\
Email & \textbf{0.683} & 0.197 ($-$0.486) & 0.683 (0.000) & -- & -- \\
Project Mgmt. & \textbf{0.652} & 0.143 ($-$0.509) & 0.652 (0.000) & -- & -- \\
SmartHome & \textbf{0.647} & 0.647 (0.000) & 0.647 (0.000) & -- & -- \\
\midrule
Macro (Full/Exact) & \textbf{0.653} & 0.369 ($-$0.284) & -- & -- & -- \\
\bottomrule
\end{tabular}
\end{table*}

\textbf{Generalized matching.}
Full exceeds Exact-only in Shopping, Calendar, Email, and Project Management, with session-bootstrap intervals excluding zero: $[0.061,0.298]$, $[0.351,0.700]$, $[0.292,0.668]$, and $[0.311,0.677]$, respectively.
Code and SmartHome are exact negative controls with identical predictions.
Thus abstraction is most useful when preferences or workflow entities vary semantically, not when normalized events already match literally.

\textbf{Phase gate.}
Removing the gate lowers Code F1 by 0.205 (95\% CI $[0.131,0.276]$) and increases FP from 23 to 59.
The Holm-corrected exact McNemar test remains significant ($p=5.89\times10^{-8}$).
The effect is repeatable across three runs, while four non-Code domains are prediction-identical with and without the gate.

\textbf{Think.}
Code w/o Think has a three-run mean F1 of 0.659 versus 0.667 for Full, and two repeats are prediction-identical to Full.
We therefore do not claim a robust aggregate contribution from the optional LLM checker.

\textbf{Cooldown.}
Removing Shopping cooldown raises ordinary F1 from 0.600 to 0.701 by increasing recall, but pushes rise from 58 to 85 and false positives from 25 to 37.
At cooldown three, 27 candidates are suppressed: 15 potential true positives and 12 potential false positives.
Cooldown therefore encodes an intervention-cost policy rather than an accuracy-only improvement.

\subsection{Stability Across Models and Protocols}

We repeat targeted Full--ablation comparisons with DeepSeek-v4-Pro, a gateway-reported GPT-5.5 endpoint, and MiniMax-M3.
Table~\ref{tab:models} reports $\Delta\mathrm{F1}=\mathrm{Full}-\mathrm{Ablation}$.
All six directions replicate, including SmartHome as a zero-effect control.
Exact score equality across several configurations reflects identical decision traces under the tested structured protocol; it should not be interpreted as independent verification of endpoint identity.

\begin{table*}[t]
\centering
\caption{Targeted mechanism effects across four model configurations. GPT-5.5 is the identity reported by the gateway and was not independently verified.}
\label{tab:models}
\small
\begin{tabular}{llrrrrc}
\toprule
Domain & Ablation & Flash & Pro & GPT-5.5 & MiniMax-M3 & Direction \\
\midrule
Code & w/o Phase & 0.205 & 0.184 & 0.184 & 0.184 & Replicated \\
Shopping & Exact-only & 0.171 & 0.171 & 0.171 & 0.171 & Replicated \\
Calendar & Exact-only & 0.539 & 0.539 & 0.539 & 0.539 & Replicated \\
Email & Exact-only & 0.486 & 0.486 & 0.486 & 0.486 & Replicated \\
Project Mgmt. & Exact-only & 0.509 & 0.509 & 0.509 & 0.509 & Replicated \\
SmartHome & Exact-only & 0.000 & 0.000 & 0.000 & 0.000 & Replicated \\
\bottomrule
\end{tabular}
\end{table*}

Code Full and Calendar Full each produce identical F1 in three observed repeats.
The shopping protocol audit additionally compares continual, split-reset, and reversed orderings; these checks preserve the principal conclusion while clarifying that the main \texttt{all} run is stateful and must not be reconstructed by summing standalone splits.

\subsection{Efficiency Composition}
\label{sec:efficiency}

Table~\ref{tab:efficiency} reports absolute \system{} telemetry under MiniMax-M3.
Across six Full runs, the framework makes 496 logical calls and observes 732,409 provider-reported tokens.
Discovery and Think account for 88 calls; optional Push formatting accounts for 408 calls (82.3\%) and 272,043 tokens (37.1\%).
Because formatting follows the structured decision, a host can replace it without changing F1.

\begin{table}[t]
\centering
\caption{Absolute Full-run telemetry under MiniMax-M3. Tokens are provider-reported observations.}
\label{tab:efficiency}
\small
\begin{tabular}{lrrr}
\toprule
Domain & Turns & Calls & Tokens \\
\midrule
Code & 448 & 105 & 112,287 \\
Shopping & 495 & 70 & 164,066 \\
Calendar & 684 & 85 & 125,510 \\
Email & 684 & 83 & 119,712 \\
Project Mgmt. & 684 & 89 & 123,218 \\
SmartHome & 440 & 64 & 87,616 \\
\midrule
Total & 3,435 & 496 & 732,409 \\
\bottomrule
\end{tabular}
\end{table}

Two optional Calendar formatting calls fail without usage payloads, so the token total is a lower bound.
Three serial Code measurements have an identical 105-call decomposition and identical decisions; mean observed tokens are 116,758 (sample SD 4,514).
We do not compare these values with matched-model baseline telemetry and therefore make no baseline-relative efficiency claim.

\subsection{Error Analysis}

Errors concentrate at the boundary between retrieving a related habit and deciding that it is actionable.
In Code-hard, 11/17 false positives occur in aligned or exception negatives, while other errors push on early edits and miss the valid wrap-up stage.
Shopping-hard similarly contains premature pushes before final decision points and false positives on preference-aligned actions.
In Calendar-all, at least 20/33 false positives are directly verified as aligned or exception cases.

These errors identify three remaining weaknesses.
First, stage cues outside Code are less explicit, so relevant habits can fire before an actionable decision.
Second, a related preference is sometimes mistaken for a violated preference.
Third, implicit semantic evidence is weaker than explicit normalized conflict keys.
The findings support the architectural separation in Figure~\ref{fig:architecture}: habit discovery can succeed while push control still fails.

\section{Related Work}

\subsection{Long-Term Memory for LLM Agents}

Generative Agents store observations, retrieve memories by relevance, recency, and importance, and synthesize higher-level reflections \citetodo{Generative Agents}.
Reflexion stores verbal feedback across trials \citetodo{Reflexion}, while MemoryBank targets persistent conversational memory and time-sensitive updating \citetodo{MemoryBank}.
MemGPT treats limited context as a hierarchy of memory tiers \citetodo{MemGPT}.
Mem0 and A-Mem study scalable extraction, consolidation, graph structure, and evolving links \citetodo{Mem0; A-Mem}.
HippoRAG uses graph structure and personalized PageRank for cross-passage integration \citetodo{HippoRAG}.
These works motivate persistent and structured memory, but their principal evaluation target is storage, retrieval, response quality, or memory organization rather than turn-level intervention timing.

\subsection{Memory Evaluation}

LoCoMo evaluates long-term conversational memory through question answering, summarization, and multimodal dialogue generation \citetodo{LoCoMo}.
LongMemEval analyzes long-term interactive memory and design choices across indexing, retrieval, and reading \citetodo{LongMemEval}.
Our evaluation differs in its output variable: each turn asks whether memory should be proactively delivered, with false positives and Push count treated as first-class outcomes.

\subsection{Proactive Memory and Agents}

PASK studies streaming latent-need detection, hybrid memory, and a complete proactive-agent system \citetodo{PASK}.
MemCog integrates proactive memory exploration with reasoning and a navigable memory store \citetodo{MemCog}.
Concurrent work, Remember When It Matters, uses a separate memory agent to decide whether to inject a reminder or remain silent during long-horizon tasks \citetodo{Remember When It Matters}.
These systems establish that proactive memory is broader than \system{}.
Our focus is the combination of (i) discovering typed habits from repeated behavior that users need not articulate, and (ii) exposing an auditable, phase- and cooldown-aware intervention controller.

\subsection{Habit and Intervention Control}

Behavioral pattern mining and interruption-aware interfaces motivate distinguishing regularity from actionability \citetodo{verified HCI work on interruption timing; workflow/phase recognition}.
\system{} operationalizes this distinction through trigger--consequence habits, phase gating, thresholding, and cooldown.
The present work evaluates offline decision labels and intervention counts; it does not replace a user study of perceived helpfulness or annoyance.

\section{Limitations}

Our six domains are adaptations designed for proactive-memory decisions, not official WorkBench, $\tau$-bench, or SmartBench protocol evaluations.
The ground truth represents designed useful-intervention cases rather than longitudinal user satisfaction.
False positives and Push count are intervention proxies; fewer pushes do not necessarily imply a better user experience.

Mechanism activation is domain dependent.
Generalized matching is inactive in Code and SmartHome, the phase gate is active primarily in Code, and the optional Think step has no robust aggregate effect in the current matrix.
The cooldown study identifies a discrete operating-point change, not a smooth or universal intervention-cost curve.

The cross-model study tests four endpoint configurations, but one identity is gateway-reported and not independently verified.
Efficiency is measured only for \system{} under MiniMax-M3, so comparisons with baseline cost or latency remain unsupported.
Finally, the system infers habits from repeated evidence but does not solve zero-history personalization, causal habit discovery, or user-facing correction and consent by itself.

\section{Conclusion}

We studied proactive memory as a decision problem rather than retrieval followed by unconditional injection.
\system{} discovers typed habits from repeated interaction evidence and separately controls whether a habit should be surfaced through generalized matching, candidate scoring, phase checks, thresholding, and cooldown.
Across six domain adaptations, it improves macro F1 and substantially reduces false and repeated interventions relative to strong memory baselines.
Ablations and error traces show that habit abstraction and delivery control solve different problems and activate under different task structures.
This separation offers a practical direction for web agents and intelligent assistants that must learn from longitudinal behavior without turning every related memory into an interruption.

\section*{Ethical Considerations}

Inferring habits creates privacy risks beyond storing explicit user statements.
The system may derive routines, preferences, schedules, emotional cues, or associations that a user did not knowingly designate as memory.
Such profiles could reveal sensitive behavior if exposed, transferred across contexts, or accessed by an unauthorized party.
A deployment should therefore require informed opt-in and provide a visible habit inspector, correction, per-habit deletion, retention limits, and a global disable control.

Proactive interventions can also be unwanted, manipulative, mistimed, or unsafe.
Threshold, phase checks, and cooldown reduce intervention frequency but do not eliminate incorrect inferences.
High-impact actions should require confirmation and application-specific policy checks.
The offline labels in this work do not establish user trust or perceived helpfulness, and the benchmark adaptations may not transfer directly to live use.

Benchmark-derived prompts were sent to external model APIs with authorization during the experiments.
A released artifact must exclude credentials and personal paths, document which benchmark fields are transmitted, and state provider retention assumptions where available.
The structured traces used for audit can themselves contain behavioral evidence and should receive the same access, retention, and deletion protections as the Habit Store.

% Bibliography intentionally omitted from this draft.
% The authors will add and verify all entries, then enable:
% \bibliographystyle{ACM-Reference-Format}
% \bibliography{references}

\end{document}
